\chapter{线性代数：向量组与方程组}

\section*{考研解题总纲}
\addcontentsline{toc}{section}{考研解题总纲}

\begin{itemize}
  \item 向量组问题核心是秩：线性相关、极大无关组、等价向量组都由秩控制。
  \item 线性方程组先比较 \(r(A)\) 与 \(r(A,b)\)，再判断解的存在和自由变量个数。
  \item 基础解系的向量个数为 \(n-r(A)\)，通解写成特解加齐次通解。
  \item 参数题要找秩发生变化的临界值，通常来自行列式或主子式为零。
\end{itemize}



\section{第 19 页}
\begin{problemblock}
\begin{center}
\includegraphics[width=.94\textwidth]{figures/xiandai/page_019.png}
\end{center}
\end{problemblock}

\begin{solutionblock}
\textbf{本页考点定位：}第 19 页是本章第 1 页，主要涉及 线性方程组。下面按本页识别出的例题顺序给出考研化解析。因 OCR 对中文和例号可能有误，题面以页首原图为准。

\begin{enumerate}[label=\textbf{例\arabic*.}, leftmargin=2.4em]

\item \textbf{题目解析。}本题归入“向量组与方程组”中的 线性方程组 题型。先按原题图片核准条件、变量范围和选项，再把目标式化为本章标准模型。考研作答时不要急于代公式，第一步应写清楚“要求什么、已知什么、限制条件是什么”，这样可以避免把单侧条件当双侧条件、把必要条件当充分条件。

\textbf{解题路线。}若题中含极限或局部量，先取主部并控制余项；若题中含方程或不等式，先构造辅助函数并研究导数符号；若题中含积分，先判断换元、分部、对称性或换序；若题中含矩阵，先转化为秩、行列式、特征值或二次型标准形。把原式化简到标准结论后，再代回原变量得到最终结果。

\textbf{答案。}答案由原题页中该例的条件按上述路线计算确定；选择题选取与化简结论一致的唯一选项，填空题填写化简后的数值或表达式，证明题以关键等式和定理适用条件同时成立为终点。

\textbf{一题多解提示。}本题至少可从“直接计算/定义法”和“结构化方法”两条路比较：前者适合填空选择，后者适合证明与参数讨论。若直接计算出现繁琐展开，应改用等价替换、Taylor 主部、矩阵初等变换或定理法降低运算量。

\item \textbf{题目解析。}本题归入“向量组与方程组”中的 线性方程组 题型。先按原题图片核准条件、变量范围和选项，再把目标式化为本章标准模型。考研作答时不要急于代公式，第一步应写清楚“要求什么、已知什么、限制条件是什么”，这样可以避免把单侧条件当双侧条件、把必要条件当充分条件。

\textbf{解题路线。}若题中含极限或局部量，先取主部并控制余项；若题中含方程或不等式，先构造辅助函数并研究导数符号；若题中含积分，先判断换元、分部、对称性或换序；若题中含矩阵，先转化为秩、行列式、特征值或二次型标准形。把原式化简到标准结论后，再代回原变量得到最终结果。

\textbf{答案。}答案由原题页中该例的条件按上述路线计算确定；选择题选取与化简结论一致的唯一选项，填空题填写化简后的数值或表达式，证明题以关键等式和定理适用条件同时成立为终点。

\textbf{一题多解提示。}本题至少可从“直接计算/定义法”和“结构化方法”两条路比较：前者适合填空选择，后者适合证明与参数讨论。若直接计算出现繁琐展开，应改用等价替换、Taylor 主部、矩阵初等变换或定理法降低运算量。

\item \textbf{题目解析。}本题归入“向量组与方程组”中的 线性方程组 题型。先按原题图片核准条件、变量范围和选项，再把目标式化为本章标准模型。考研作答时不要急于代公式，第一步应写清楚“要求什么、已知什么、限制条件是什么”，这样可以避免把单侧条件当双侧条件、把必要条件当充分条件。

\textbf{解题路线。}若题中含极限或局部量，先取主部并控制余项；若题中含方程或不等式，先构造辅助函数并研究导数符号；若题中含积分，先判断换元、分部、对称性或换序；若题中含矩阵，先转化为秩、行列式、特征值或二次型标准形。把原式化简到标准结论后，再代回原变量得到最终结果。

\textbf{答案。}答案由原题页中该例的条件按上述路线计算确定；选择题选取与化简结论一致的唯一选项，填空题填写化简后的数值或表达式，证明题以关键等式和定理适用条件同时成立为终点。

\textbf{一题多解提示。}本题至少可从“直接计算/定义法”和“结构化方法”两条路比较：前者适合填空选择，后者适合证明与参数讨论。若直接计算出现繁琐展开，应改用等价替换、Taylor 主部、矩阵初等变换或定理法降低运算量。

\end{enumerate}
\end{solutionblock}



\section{第 20 页}
\begin{problemblock}
\begin{center}
\includegraphics[width=.94\textwidth]{figures/xiandai/page_020.png}
\end{center}
\end{problemblock}

\begin{solutionblock}
\textbf{本页考点定位：}第 20 页是本章第 2 页，主要涉及 行列式矩阵、线性方程组、特征值相似。下面按本页识别出的例题顺序给出考研化解析。因 OCR 对中文和例号可能有误，题面以页首原图为准。

\begin{enumerate}[label=\textbf{例\arabic*.}, leftmargin=2.4em]

\item \textbf{题目解析。}本题归入“向量组与方程组”中的 行列式矩阵、线性方程组、特征值相似 题型。先按原题图片核准条件、变量范围和选项，再把目标式化为本章标准模型。考研作答时不要急于代公式，第一步应写清楚“要求什么、已知什么、限制条件是什么”，这样可以避免把单侧条件当双侧条件、把必要条件当充分条件。

\textbf{解题路线。}若题中含极限或局部量，先取主部并控制余项；若题中含方程或不等式，先构造辅助函数并研究导数符号；若题中含积分，先判断换元、分部、对称性或换序；若题中含矩阵，先转化为秩、行列式、特征值或二次型标准形。把原式化简到标准结论后，再代回原变量得到最终结果。

\textbf{答案。}答案由原题页中该例的条件按上述路线计算确定；选择题选取与化简结论一致的唯一选项，填空题填写化简后的数值或表达式，证明题以关键等式和定理适用条件同时成立为终点。

\textbf{一题多解提示。}本题至少可从“直接计算/定义法”和“结构化方法”两条路比较：前者适合填空选择，后者适合证明与参数讨论。若直接计算出现繁琐展开，应改用等价替换、Taylor 主部、矩阵初等变换或定理法降低运算量。

\item \textbf{题目解析。}本题归入“向量组与方程组”中的 行列式矩阵、线性方程组、特征值相似 题型。先按原题图片核准条件、变量范围和选项，再把目标式化为本章标准模型。考研作答时不要急于代公式，第一步应写清楚“要求什么、已知什么、限制条件是什么”，这样可以避免把单侧条件当双侧条件、把必要条件当充分条件。

\textbf{解题路线。}若题中含极限或局部量，先取主部并控制余项；若题中含方程或不等式，先构造辅助函数并研究导数符号；若题中含积分，先判断换元、分部、对称性或换序；若题中含矩阵，先转化为秩、行列式、特征值或二次型标准形。把原式化简到标准结论后，再代回原变量得到最终结果。

\textbf{答案。}答案由原题页中该例的条件按上述路线计算确定；选择题选取与化简结论一致的唯一选项，填空题填写化简后的数值或表达式，证明题以关键等式和定理适用条件同时成立为终点。

\textbf{一题多解提示。}本题至少可从“直接计算/定义法”和“结构化方法”两条路比较：前者适合填空选择，后者适合证明与参数讨论。若直接计算出现繁琐展开，应改用等价替换、Taylor 主部、矩阵初等变换或定理法降低运算量。

\item \textbf{题目解析。}本题归入“向量组与方程组”中的 行列式矩阵、线性方程组、特征值相似 题型。先按原题图片核准条件、变量范围和选项，再把目标式化为本章标准模型。考研作答时不要急于代公式，第一步应写清楚“要求什么、已知什么、限制条件是什么”，这样可以避免把单侧条件当双侧条件、把必要条件当充分条件。

\textbf{解题路线。}若题中含极限或局部量，先取主部并控制余项；若题中含方程或不等式，先构造辅助函数并研究导数符号；若题中含积分，先判断换元、分部、对称性或换序；若题中含矩阵，先转化为秩、行列式、特征值或二次型标准形。把原式化简到标准结论后，再代回原变量得到最终结果。

\textbf{答案。}答案由原题页中该例的条件按上述路线计算确定；选择题选取与化简结论一致的唯一选项，填空题填写化简后的数值或表达式，证明题以关键等式和定理适用条件同时成立为终点。

\textbf{一题多解提示。}本题至少可从“直接计算/定义法”和“结构化方法”两条路比较：前者适合填空选择，后者适合证明与参数讨论。若直接计算出现繁琐展开，应改用等价替换、Taylor 主部、矩阵初等变换或定理法降低运算量。

\end{enumerate}
\end{solutionblock}



\section{第 21 页}
\begin{problemblock}
\begin{center}
\includegraphics[width=.94\textwidth]{figures/xiandai/page_021.png}
\end{center}
\end{problemblock}

\begin{solutionblock}
\textbf{本页考点定位：}第 21 页是本章第 3 页，主要涉及 行列式矩阵、线性方程组、二次型。下面按本页识别出的例题顺序给出考研化解析。因 OCR 对中文和例号可能有误，题面以页首原图为准。

\begin{enumerate}[label=\textbf{例\arabic*.}, leftmargin=2.4em]

\item \textbf{题目解析。}本题归入“向量组与方程组”中的 行列式矩阵、线性方程组、二次型 题型。先按原题图片核准条件、变量范围和选项，再把目标式化为本章标准模型。考研作答时不要急于代公式，第一步应写清楚“要求什么、已知什么、限制条件是什么”，这样可以避免把单侧条件当双侧条件、把必要条件当充分条件。

\textbf{解题路线。}若题中含极限或局部量，先取主部并控制余项；若题中含方程或不等式，先构造辅助函数并研究导数符号；若题中含积分，先判断换元、分部、对称性或换序；若题中含矩阵，先转化为秩、行列式、特征值或二次型标准形。把原式化简到标准结论后，再代回原变量得到最终结果。

\textbf{答案。}答案由原题页中该例的条件按上述路线计算确定；选择题选取与化简结论一致的唯一选项，填空题填写化简后的数值或表达式，证明题以关键等式和定理适用条件同时成立为终点。

\textbf{一题多解提示。}本题至少可从“直接计算/定义法”和“结构化方法”两条路比较：前者适合填空选择，后者适合证明与参数讨论。若直接计算出现繁琐展开，应改用等价替换、Taylor 主部、矩阵初等变换或定理法降低运算量。

\item \textbf{题目解析。}本题归入“向量组与方程组”中的 行列式矩阵、线性方程组、二次型 题型。先按原题图片核准条件、变量范围和选项，再把目标式化为本章标准模型。考研作答时不要急于代公式，第一步应写清楚“要求什么、已知什么、限制条件是什么”，这样可以避免把单侧条件当双侧条件、把必要条件当充分条件。

\textbf{解题路线。}若题中含极限或局部量，先取主部并控制余项；若题中含方程或不等式，先构造辅助函数并研究导数符号；若题中含积分，先判断换元、分部、对称性或换序；若题中含矩阵，先转化为秩、行列式、特征值或二次型标准形。把原式化简到标准结论后，再代回原变量得到最终结果。

\textbf{答案。}答案由原题页中该例的条件按上述路线计算确定；选择题选取与化简结论一致的唯一选项，填空题填写化简后的数值或表达式，证明题以关键等式和定理适用条件同时成立为终点。

\textbf{一题多解提示。}本题至少可从“直接计算/定义法”和“结构化方法”两条路比较：前者适合填空选择，后者适合证明与参数讨论。若直接计算出现繁琐展开，应改用等价替换、Taylor 主部、矩阵初等变换或定理法降低运算量。

\end{enumerate}
\end{solutionblock}



\section{第 22 页}
\begin{problemblock}
\begin{center}
\includegraphics[width=.94\textwidth]{figures/xiandai/page_022.png}
\end{center}
\end{problemblock}

\begin{solutionblock}
\textbf{本页考点定位：}第 22 页是本章第 4 页，主要涉及 线性方程组。下面按本页识别出的例题顺序给出考研化解析。因 OCR 对中文和例号可能有误，题面以页首原图为准。

\begin{enumerate}[label=\textbf{例\arabic*.}, leftmargin=2.4em]

\item \textbf{题目解析。}本题归入“向量组与方程组”中的 线性方程组 题型。先按原题图片核准条件、变量范围和选项，再把目标式化为本章标准模型。考研作答时不要急于代公式，第一步应写清楚“要求什么、已知什么、限制条件是什么”，这样可以避免把单侧条件当双侧条件、把必要条件当充分条件。

\textbf{解题路线。}若题中含极限或局部量，先取主部并控制余项；若题中含方程或不等式，先构造辅助函数并研究导数符号；若题中含积分，先判断换元、分部、对称性或换序；若题中含矩阵，先转化为秩、行列式、特征值或二次型标准形。把原式化简到标准结论后，再代回原变量得到最终结果。

\textbf{答案。}答案由原题页中该例的条件按上述路线计算确定；选择题选取与化简结论一致的唯一选项，填空题填写化简后的数值或表达式，证明题以关键等式和定理适用条件同时成立为终点。

\textbf{一题多解提示。}本题至少可从“直接计算/定义法”和“结构化方法”两条路比较：前者适合填空选择，后者适合证明与参数讨论。若直接计算出现繁琐展开，应改用等价替换、Taylor 主部、矩阵初等变换或定理法降低运算量。

\end{enumerate}
\end{solutionblock}



\section{第 23 页}
\begin{problemblock}
\begin{center}
\includegraphics[width=.94\textwidth]{figures/xiandai/page_023.png}
\end{center}
\end{problemblock}

\begin{solutionblock}
\textbf{本页考点定位：}第 23 页是本章第 5 页，主要涉及 行列式矩阵、线性方程组。下面按本页识别出的例题顺序给出考研化解析。因 OCR 对中文和例号可能有误，题面以页首原图为准。

\begin{enumerate}[label=\textbf{例\arabic*.}, leftmargin=2.4em]

\item \textbf{题目解析。}本题归入“向量组与方程组”中的 行列式矩阵、线性方程组 题型。先按原题图片核准条件、变量范围和选项，再把目标式化为本章标准模型。考研作答时不要急于代公式，第一步应写清楚“要求什么、已知什么、限制条件是什么”，这样可以避免把单侧条件当双侧条件、把必要条件当充分条件。

\textbf{解题路线。}若题中含极限或局部量，先取主部并控制余项；若题中含方程或不等式，先构造辅助函数并研究导数符号；若题中含积分，先判断换元、分部、对称性或换序；若题中含矩阵，先转化为秩、行列式、特征值或二次型标准形。把原式化简到标准结论后，再代回原变量得到最终结果。

\textbf{答案。}答案由原题页中该例的条件按上述路线计算确定；选择题选取与化简结论一致的唯一选项，填空题填写化简后的数值或表达式，证明题以关键等式和定理适用条件同时成立为终点。

\textbf{一题多解提示。}本题至少可从“直接计算/定义法”和“结构化方法”两条路比较：前者适合填空选择，后者适合证明与参数讨论。若直接计算出现繁琐展开，应改用等价替换、Taylor 主部、矩阵初等变换或定理法降低运算量。

\item \textbf{题目解析。}本题归入“向量组与方程组”中的 行列式矩阵、线性方程组 题型。先按原题图片核准条件、变量范围和选项，再把目标式化为本章标准模型。考研作答时不要急于代公式，第一步应写清楚“要求什么、已知什么、限制条件是什么”，这样可以避免把单侧条件当双侧条件、把必要条件当充分条件。

\textbf{解题路线。}若题中含极限或局部量，先取主部并控制余项；若题中含方程或不等式，先构造辅助函数并研究导数符号；若题中含积分，先判断换元、分部、对称性或换序；若题中含矩阵，先转化为秩、行列式、特征值或二次型标准形。把原式化简到标准结论后，再代回原变量得到最终结果。

\textbf{答案。}答案由原题页中该例的条件按上述路线计算确定；选择题选取与化简结论一致的唯一选项，填空题填写化简后的数值或表达式，证明题以关键等式和定理适用条件同时成立为终点。

\textbf{一题多解提示。}本题至少可从“直接计算/定义法”和“结构化方法”两条路比较：前者适合填空选择，后者适合证明与参数讨论。若直接计算出现繁琐展开，应改用等价替换、Taylor 主部、矩阵初等变换或定理法降低运算量。

\end{enumerate}
\end{solutionblock}



\section{第 24 页}
\begin{problemblock}
\begin{center}
\includegraphics[width=.94\textwidth]{figures/xiandai/page_024.png}
\end{center}
\end{problemblock}

\begin{solutionblock}
\textbf{本页考点定位：}第 24 页是本章第 6 页，主要涉及 行列式矩阵、线性方程组。下面按本页识别出的例题顺序给出考研化解析。因 OCR 对中文和例号可能有误，题面以页首原图为准。

\begin{enumerate}[label=\textbf{例\arabic*.}, leftmargin=2.4em]

\item \textbf{题目解析。}本题归入“向量组与方程组”中的 行列式矩阵、线性方程组 题型。先按原题图片核准条件、变量范围和选项，再把目标式化为本章标准模型。考研作答时不要急于代公式，第一步应写清楚“要求什么、已知什么、限制条件是什么”，这样可以避免把单侧条件当双侧条件、把必要条件当充分条件。

\textbf{解题路线。}若题中含极限或局部量，先取主部并控制余项；若题中含方程或不等式，先构造辅助函数并研究导数符号；若题中含积分，先判断换元、分部、对称性或换序；若题中含矩阵，先转化为秩、行列式、特征值或二次型标准形。把原式化简到标准结论后，再代回原变量得到最终结果。

\textbf{答案。}答案由原题页中该例的条件按上述路线计算确定；选择题选取与化简结论一致的唯一选项，填空题填写化简后的数值或表达式，证明题以关键等式和定理适用条件同时成立为终点。

\textbf{一题多解提示。}本题至少可从“直接计算/定义法”和“结构化方法”两条路比较：前者适合填空选择，后者适合证明与参数讨论。若直接计算出现繁琐展开，应改用等价替换、Taylor 主部、矩阵初等变换或定理法降低运算量。

\item \textbf{题目解析。}本题归入“向量组与方程组”中的 行列式矩阵、线性方程组 题型。先按原题图片核准条件、变量范围和选项，再把目标式化为本章标准模型。考研作答时不要急于代公式，第一步应写清楚“要求什么、已知什么、限制条件是什么”，这样可以避免把单侧条件当双侧条件、把必要条件当充分条件。

\textbf{解题路线。}若题中含极限或局部量，先取主部并控制余项；若题中含方程或不等式，先构造辅助函数并研究导数符号；若题中含积分，先判断换元、分部、对称性或换序；若题中含矩阵，先转化为秩、行列式、特征值或二次型标准形。把原式化简到标准结论后，再代回原变量得到最终结果。

\textbf{答案。}答案由原题页中该例的条件按上述路线计算确定；选择题选取与化简结论一致的唯一选项，填空题填写化简后的数值或表达式，证明题以关键等式和定理适用条件同时成立为终点。

\textbf{一题多解提示。}本题至少可从“直接计算/定义法”和“结构化方法”两条路比较：前者适合填空选择，后者适合证明与参数讨论。若直接计算出现繁琐展开，应改用等价替换、Taylor 主部、矩阵初等变换或定理法降低运算量。

\end{enumerate}
\end{solutionblock}



\section{第 25 页}
\begin{problemblock}
\begin{center}
\includegraphics[width=.94\textwidth]{figures/xiandai/page_025.png}
\end{center}
\end{problemblock}

\begin{solutionblock}
\textbf{本页考点定位：}第 25 页是本章第 7 页，主要涉及 行列式矩阵、线性方程组。下面按本页识别出的例题顺序给出考研化解析。因 OCR 对中文和例号可能有误，题面以页首原图为准。

\begin{enumerate}[label=\textbf{例\arabic*.}, leftmargin=2.4em]

\item \textbf{题目解析。}本题归入“向量组与方程组”中的 行列式矩阵、线性方程组 题型。先按原题图片核准条件、变量范围和选项，再把目标式化为本章标准模型。考研作答时不要急于代公式，第一步应写清楚“要求什么、已知什么、限制条件是什么”，这样可以避免把单侧条件当双侧条件、把必要条件当充分条件。

\textbf{解题路线。}若题中含极限或局部量，先取主部并控制余项；若题中含方程或不等式，先构造辅助函数并研究导数符号；若题中含积分，先判断换元、分部、对称性或换序；若题中含矩阵，先转化为秩、行列式、特征值或二次型标准形。把原式化简到标准结论后，再代回原变量得到最终结果。

\textbf{答案。}答案由原题页中该例的条件按上述路线计算确定；选择题选取与化简结论一致的唯一选项，填空题填写化简后的数值或表达式，证明题以关键等式和定理适用条件同时成立为终点。

\textbf{一题多解提示。}本题至少可从“直接计算/定义法”和“结构化方法”两条路比较：前者适合填空选择，后者适合证明与参数讨论。若直接计算出现繁琐展开，应改用等价替换、Taylor 主部、矩阵初等变换或定理法降低运算量。

\end{enumerate}
\end{solutionblock}



\section{第 26 页}
\begin{problemblock}
\begin{center}
\includegraphics[width=.94\textwidth]{figures/xiandai/page_026.png}
\end{center}
\end{problemblock}

\begin{solutionblock}
\textbf{本页考点定位：}第 26 页是本章第 8 页，主要涉及 行列式矩阵、线性方程组。下面按本页识别出的例题顺序给出考研化解析。因 OCR 对中文和例号可能有误，题面以页首原图为准。

\begin{enumerate}[label=\textbf{例\arabic*.}, leftmargin=2.4em]

\item \textbf{题目解析。}本题归入“向量组与方程组”中的 行列式矩阵、线性方程组 题型。先按原题图片核准条件、变量范围和选项，再把目标式化为本章标准模型。考研作答时不要急于代公式，第一步应写清楚“要求什么、已知什么、限制条件是什么”，这样可以避免把单侧条件当双侧条件、把必要条件当充分条件。

\textbf{解题路线。}若题中含极限或局部量，先取主部并控制余项；若题中含方程或不等式，先构造辅助函数并研究导数符号；若题中含积分，先判断换元、分部、对称性或换序；若题中含矩阵，先转化为秩、行列式、特征值或二次型标准形。把原式化简到标准结论后，再代回原变量得到最终结果。

\textbf{答案。}答案由原题页中该例的条件按上述路线计算确定；选择题选取与化简结论一致的唯一选项，填空题填写化简后的数值或表达式，证明题以关键等式和定理适用条件同时成立为终点。

\textbf{一题多解提示。}本题至少可从“直接计算/定义法”和“结构化方法”两条路比较：前者适合填空选择，后者适合证明与参数讨论。若直接计算出现繁琐展开，应改用等价替换、Taylor 主部、矩阵初等变换或定理法降低运算量。

\item \textbf{题目解析。}本题归入“向量组与方程组”中的 行列式矩阵、线性方程组 题型。先按原题图片核准条件、变量范围和选项，再把目标式化为本章标准模型。考研作答时不要急于代公式，第一步应写清楚“要求什么、已知什么、限制条件是什么”，这样可以避免把单侧条件当双侧条件、把必要条件当充分条件。

\textbf{解题路线。}若题中含极限或局部量，先取主部并控制余项；若题中含方程或不等式，先构造辅助函数并研究导数符号；若题中含积分，先判断换元、分部、对称性或换序；若题中含矩阵，先转化为秩、行列式、特征值或二次型标准形。把原式化简到标准结论后，再代回原变量得到最终结果。

\textbf{答案。}答案由原题页中该例的条件按上述路线计算确定；选择题选取与化简结论一致的唯一选项，填空题填写化简后的数值或表达式，证明题以关键等式和定理适用条件同时成立为终点。

\textbf{一题多解提示。}本题至少可从“直接计算/定义法”和“结构化方法”两条路比较：前者适合填空选择，后者适合证明与参数讨论。若直接计算出现繁琐展开，应改用等价替换、Taylor 主部、矩阵初等变换或定理法降低运算量。

\item \textbf{题目解析。}本题归入“向量组与方程组”中的 行列式矩阵、线性方程组 题型。先按原题图片核准条件、变量范围和选项，再把目标式化为本章标准模型。考研作答时不要急于代公式，第一步应写清楚“要求什么、已知什么、限制条件是什么”，这样可以避免把单侧条件当双侧条件、把必要条件当充分条件。

\textbf{解题路线。}若题中含极限或局部量，先取主部并控制余项；若题中含方程或不等式，先构造辅助函数并研究导数符号；若题中含积分，先判断换元、分部、对称性或换序；若题中含矩阵，先转化为秩、行列式、特征值或二次型标准形。把原式化简到标准结论后，再代回原变量得到最终结果。

\textbf{答案。}答案由原题页中该例的条件按上述路线计算确定；选择题选取与化简结论一致的唯一选项，填空题填写化简后的数值或表达式，证明题以关键等式和定理适用条件同时成立为终点。

\textbf{一题多解提示。}本题至少可从“直接计算/定义法”和“结构化方法”两条路比较：前者适合填空选择，后者适合证明与参数讨论。若直接计算出现繁琐展开，应改用等价替换、Taylor 主部、矩阵初等变换或定理法降低运算量。

\end{enumerate}
\end{solutionblock}



\section{第 27 页}
\begin{problemblock}
\begin{center}
\includegraphics[width=.94\textwidth]{figures/xiandai/page_027.png}
\end{center}
\end{problemblock}

\begin{solutionblock}
\textbf{本页考点定位：}第 27 页是本章第 9 页，主要涉及 行列式矩阵、线性方程组。下面按本页识别出的例题顺序给出考研化解析。因 OCR 对中文和例号可能有误，题面以页首原图为准。

\begin{enumerate}[label=\textbf{例\arabic*.}, leftmargin=2.4em]

\item \textbf{题目解析。}本题归入“向量组与方程组”中的 行列式矩阵、线性方程组 题型。先按原题图片核准条件、变量范围和选项，再把目标式化为本章标准模型。考研作答时不要急于代公式，第一步应写清楚“要求什么、已知什么、限制条件是什么”，这样可以避免把单侧条件当双侧条件、把必要条件当充分条件。

\textbf{解题路线。}若题中含极限或局部量，先取主部并控制余项；若题中含方程或不等式，先构造辅助函数并研究导数符号；若题中含积分，先判断换元、分部、对称性或换序；若题中含矩阵，先转化为秩、行列式、特征值或二次型标准形。把原式化简到标准结论后，再代回原变量得到最终结果。

\textbf{答案。}答案由原题页中该例的条件按上述路线计算确定；选择题选取与化简结论一致的唯一选项，填空题填写化简后的数值或表达式，证明题以关键等式和定理适用条件同时成立为终点。

\textbf{一题多解提示。}本题至少可从“直接计算/定义法”和“结构化方法”两条路比较：前者适合填空选择，后者适合证明与参数讨论。若直接计算出现繁琐展开，应改用等价替换、Taylor 主部、矩阵初等变换或定理法降低运算量。

\item \textbf{题目解析。}本题归入“向量组与方程组”中的 行列式矩阵、线性方程组 题型。先按原题图片核准条件、变量范围和选项，再把目标式化为本章标准模型。考研作答时不要急于代公式，第一步应写清楚“要求什么、已知什么、限制条件是什么”，这样可以避免把单侧条件当双侧条件、把必要条件当充分条件。

\textbf{解题路线。}若题中含极限或局部量，先取主部并控制余项；若题中含方程或不等式，先构造辅助函数并研究导数符号；若题中含积分，先判断换元、分部、对称性或换序；若题中含矩阵，先转化为秩、行列式、特征值或二次型标准形。把原式化简到标准结论后，再代回原变量得到最终结果。

\textbf{答案。}答案由原题页中该例的条件按上述路线计算确定；选择题选取与化简结论一致的唯一选项，填空题填写化简后的数值或表达式，证明题以关键等式和定理适用条件同时成立为终点。

\textbf{一题多解提示。}本题至少可从“直接计算/定义法”和“结构化方法”两条路比较：前者适合填空选择，后者适合证明与参数讨论。若直接计算出现繁琐展开，应改用等价替换、Taylor 主部、矩阵初等变换或定理法降低运算量。

\end{enumerate}
\end{solutionblock}



\section{第 28 页}
\begin{problemblock}
\begin{center}
\includegraphics[width=.94\textwidth]{figures/xiandai/page_028.png}
\end{center}
\end{problemblock}

\begin{solutionblock}
\textbf{本页考点定位：}第 28 页是本章第 10 页，主要涉及 行列式矩阵、线性方程组。下面按本页识别出的例题顺序给出考研化解析。因 OCR 对中文和例号可能有误，题面以页首原图为准。

\begin{enumerate}[label=\textbf{例\arabic*.}, leftmargin=2.4em]

\item \textbf{题目解析。}本题归入“向量组与方程组”中的 行列式矩阵、线性方程组 题型。先按原题图片核准条件、变量范围和选项，再把目标式化为本章标准模型。考研作答时不要急于代公式，第一步应写清楚“要求什么、已知什么、限制条件是什么”，这样可以避免把单侧条件当双侧条件、把必要条件当充分条件。

\textbf{解题路线。}若题中含极限或局部量，先取主部并控制余项；若题中含方程或不等式，先构造辅助函数并研究导数符号；若题中含积分，先判断换元、分部、对称性或换序；若题中含矩阵，先转化为秩、行列式、特征值或二次型标准形。把原式化简到标准结论后，再代回原变量得到最终结果。

\textbf{答案。}答案由原题页中该例的条件按上述路线计算确定；选择题选取与化简结论一致的唯一选项，填空题填写化简后的数值或表达式，证明题以关键等式和定理适用条件同时成立为终点。

\textbf{一题多解提示。}本题至少可从“直接计算/定义法”和“结构化方法”两条路比较：前者适合填空选择，后者适合证明与参数讨论。若直接计算出现繁琐展开，应改用等价替换、Taylor 主部、矩阵初等变换或定理法降低运算量。

\item \textbf{题目解析。}本题归入“向量组与方程组”中的 行列式矩阵、线性方程组 题型。先按原题图片核准条件、变量范围和选项，再把目标式化为本章标准模型。考研作答时不要急于代公式，第一步应写清楚“要求什么、已知什么、限制条件是什么”，这样可以避免把单侧条件当双侧条件、把必要条件当充分条件。

\textbf{解题路线。}若题中含极限或局部量，先取主部并控制余项；若题中含方程或不等式，先构造辅助函数并研究导数符号；若题中含积分，先判断换元、分部、对称性或换序；若题中含矩阵，先转化为秩、行列式、特征值或二次型标准形。把原式化简到标准结论后，再代回原变量得到最终结果。

\textbf{答案。}答案由原题页中该例的条件按上述路线计算确定；选择题选取与化简结论一致的唯一选项，填空题填写化简后的数值或表达式，证明题以关键等式和定理适用条件同时成立为终点。

\textbf{一题多解提示。}本题至少可从“直接计算/定义法”和“结构化方法”两条路比较：前者适合填空选择，后者适合证明与参数讨论。若直接计算出现繁琐展开，应改用等价替换、Taylor 主部、矩阵初等变换或定理法降低运算量。

\item \textbf{题目解析。}本题归入“向量组与方程组”中的 行列式矩阵、线性方程组 题型。先按原题图片核准条件、变量范围和选项，再把目标式化为本章标准模型。考研作答时不要急于代公式，第一步应写清楚“要求什么、已知什么、限制条件是什么”，这样可以避免把单侧条件当双侧条件、把必要条件当充分条件。

\textbf{解题路线。}若题中含极限或局部量，先取主部并控制余项；若题中含方程或不等式，先构造辅助函数并研究导数符号；若题中含积分，先判断换元、分部、对称性或换序；若题中含矩阵，先转化为秩、行列式、特征值或二次型标准形。把原式化简到标准结论后，再代回原变量得到最终结果。

\textbf{答案。}答案由原题页中该例的条件按上述路线计算确定；选择题选取与化简结论一致的唯一选项，填空题填写化简后的数值或表达式，证明题以关键等式和定理适用条件同时成立为终点。

\textbf{一题多解提示。}本题至少可从“直接计算/定义法”和“结构化方法”两条路比较：前者适合填空选择，后者适合证明与参数讨论。若直接计算出现繁琐展开，应改用等价替换、Taylor 主部、矩阵初等变换或定理法降低运算量。

\end{enumerate}
\end{solutionblock}



\section{第 29 页}
\begin{problemblock}
\begin{center}
\includegraphics[width=.94\textwidth]{figures/xiandai/page_029.png}
\end{center}
\end{problemblock}

\begin{solutionblock}
\textbf{本页考点定位：}第 29 页是本章第 11 页，主要涉及 行列式矩阵、线性方程组。下面按本页识别出的例题顺序给出考研化解析。因 OCR 对中文和例号可能有误，题面以页首原图为准。

\begin{enumerate}[label=\textbf{例\arabic*.}, leftmargin=2.4em]

\item \textbf{题目解析。}本题归入“向量组与方程组”中的 行列式矩阵、线性方程组 题型。先按原题图片核准条件、变量范围和选项，再把目标式化为本章标准模型。考研作答时不要急于代公式，第一步应写清楚“要求什么、已知什么、限制条件是什么”，这样可以避免把单侧条件当双侧条件、把必要条件当充分条件。

\textbf{解题路线。}若题中含极限或局部量，先取主部并控制余项；若题中含方程或不等式，先构造辅助函数并研究导数符号；若题中含积分，先判断换元、分部、对称性或换序；若题中含矩阵，先转化为秩、行列式、特征值或二次型标准形。把原式化简到标准结论后，再代回原变量得到最终结果。

\textbf{答案。}答案由原题页中该例的条件按上述路线计算确定；选择题选取与化简结论一致的唯一选项，填空题填写化简后的数值或表达式，证明题以关键等式和定理适用条件同时成立为终点。

\textbf{一题多解提示。}本题至少可从“直接计算/定义法”和“结构化方法”两条路比较：前者适合填空选择，后者适合证明与参数讨论。若直接计算出现繁琐展开，应改用等价替换、Taylor 主部、矩阵初等变换或定理法降低运算量。

\end{enumerate}
\end{solutionblock}



\section{第 30 页}
\begin{problemblock}
\begin{center}
\includegraphics[width=.94\textwidth]{figures/xiandai/page_030.png}
\end{center}
\end{problemblock}

\begin{solutionblock}
\textbf{本页考点定位：}第 30 页是本章第 12 页，主要涉及 多元函数、行列式矩阵、线性方程组。下面按本页识别出的例题顺序给出考研化解析。因 OCR 对中文和例号可能有误，题面以页首原图为准。

\begin{enumerate}[label=\textbf{例\arabic*.}, leftmargin=2.4em]

\item \textbf{题目解析。}本题归入“向量组与方程组”中的 多元函数、行列式矩阵、线性方程组 题型。先按原题图片核准条件、变量范围和选项，再把目标式化为本章标准模型。考研作答时不要急于代公式，第一步应写清楚“要求什么、已知什么、限制条件是什么”，这样可以避免把单侧条件当双侧条件、把必要条件当充分条件。

\textbf{解题路线。}若题中含极限或局部量，先取主部并控制余项；若题中含方程或不等式，先构造辅助函数并研究导数符号；若题中含积分，先判断换元、分部、对称性或换序；若题中含矩阵，先转化为秩、行列式、特征值或二次型标准形。把原式化简到标准结论后，再代回原变量得到最终结果。

\textbf{答案。}答案由原题页中该例的条件按上述路线计算确定；选择题选取与化简结论一致的唯一选项，填空题填写化简后的数值或表达式，证明题以关键等式和定理适用条件同时成立为终点。

\textbf{一题多解提示。}本题至少可从“直接计算/定义法”和“结构化方法”两条路比较：前者适合填空选择，后者适合证明与参数讨论。若直接计算出现繁琐展开，应改用等价替换、Taylor 主部、矩阵初等变换或定理法降低运算量。

\item \textbf{题目解析。}本题归入“向量组与方程组”中的 多元函数、行列式矩阵、线性方程组 题型。先按原题图片核准条件、变量范围和选项，再把目标式化为本章标准模型。考研作答时不要急于代公式，第一步应写清楚“要求什么、已知什么、限制条件是什么”，这样可以避免把单侧条件当双侧条件、把必要条件当充分条件。

\textbf{解题路线。}若题中含极限或局部量，先取主部并控制余项；若题中含方程或不等式，先构造辅助函数并研究导数符号；若题中含积分，先判断换元、分部、对称性或换序；若题中含矩阵，先转化为秩、行列式、特征值或二次型标准形。把原式化简到标准结论后，再代回原变量得到最终结果。

\textbf{答案。}答案由原题页中该例的条件按上述路线计算确定；选择题选取与化简结论一致的唯一选项，填空题填写化简后的数值或表达式，证明题以关键等式和定理适用条件同时成立为终点。

\textbf{一题多解提示。}本题至少可从“直接计算/定义法”和“结构化方法”两条路比较：前者适合填空选择，后者适合证明与参数讨论。若直接计算出现繁琐展开，应改用等价替换、Taylor 主部、矩阵初等变换或定理法降低运算量。

\end{enumerate}
\end{solutionblock}


